\begin{proposition}[Active nonterminal branch stays in the unmodified $B$-region, conditional on the raw active law]
Fix odd $m$ and the best-seed branch for the unresolved channel $R1\to H_{L1}$. Let
\[
\Theta:=q+s+v+\lambda \pmod m.
\]
For the seed $L=[2,2,1]$, $R=[1,4,4]$ with $(w_0,s_0)=(0,0)$, the six patched classes are exactly
\[
R1 = \{\Theta=1,\ q=m-1,\ w=0,\ s\neq 0\},
\quad
L1 = \{\Theta=1,\ q=m-1,\ w=m-1,\ s\neq 0\},
\]
\[
R2 = \{\Theta=2,\ q=0,\ w=0,\ s\neq 0\},
\quad
L2 = \{\Theta=2,\ q=m-1,\ w=0,\ s\neq 1\},
\]
\[
R3 = \{\Theta=3,\ q=0,\ w=0,\ s\neq 1\},
\quad
L3 = \{\Theta=3,\ q=m-1,\ w=1,\ s\neq 2\}.
\]
Assume the active nonterminal branch starts at
\[
(m-1,1,2)
\]
and follows the raw current-coordinate law
\[
(q,w,0)\mapsto(q+1,w,1),
\qquad
(q,w,1)\mapsto(q,w,2),
\]
\[
(q,w,2)\mapsto(q,w+\mathbf 1_{\{q=m-1\}},3),
\qquad
(q,w,\Theta)\mapsto(q,w,\Theta+1)
\quad (3\le \Theta\le m-1),
\]
until the first exit states
\[
(m-1,m-2,1) \quad \text{or} \quad (m-2,m-1,1).
\]
Then every nonterminal active state is $B$-labeled.
\end{proposition}

\begin{proof}[Proof sketch]
Along the active branch, $w$ starts at $1$ and can only increase, never decrease. Since the branch exits at $(m-1,m-2,1)$ or $(m-2,m-1,1)$ before $w$ can wrap, one has $w\neq 0$ on every nonterminal state. This excludes $R1,R2,R3,L2$.

The value $w=1$ occurs only at the entry state $(m-1,1,2)$: the very next step sends it to $(m-1,2,3)$, and $w$ never decreases afterward. Hence $L3$ is never visited.

Finally, $w=m-1$ can only be created from a phase-$2$ state $(m-1,m-2,2)$, whose successor is $(m-1,m-1,3)$. While $q=m-1$, the phase then runs through $3,4,\dots,m-1,0$, and the next step changes $q$ to $0$ at phase $1$. Therefore $(m-1,m-1,1)$ never occurs, so $L1$ is also excluded.

Hence none of the six patched classes is visited, and every nonterminal active state is $B$-labeled.\qedhere
\end{proof}
